\chapter{Conclusion and Future Work}
\label{Chapter6}

\section{Summary}

This thesis set out to evaluate LLM-based Android malware detection under a
common framework, treating detection quality and cost-efficiency as
co-equal evaluation axes rather than reporting cost as an afterthought to
accuracy (Chapter~\ref{Chapter4}). The first phase was a Systematic
Literature Review of six primary studies (Chapter~\ref{Chapter3}), which
mapped the field into three recurring strategies (Retrieval-Augmented
Generation, Semantic Representation, and Fine-tuning/Direct
Classification), synthesized their reported effectiveness and cost
profiles without over-claiming direct numerical comparability, and
identified concrete gaps: inconsistent cost reporting, a shared reliance
on static analysis alone, and reliability limits from hallucination and
prompt sensitivity. The second phase built and iteratively hardened one
concrete candidate on top of the LAMD framework (Chapter~\ref{Chapter4}),
progressing from a faithful baseline reimplementation through role-based
early exit, a semantic cache, and multi-level SOURCE/SINK connection
verification, with each optimization's own accuracy risk found and fixed
before being accepted (Table~\ref{tab:design-history}). The measured
results (Chapter~\ref{Chapter5}) show the two cost-reduction mechanisms
working as designed -- up to 85\% fewer LLM calls and 72\% fewer tokens
from the cache alone, 51--97\% fewer calls from connection-verified early
exit -- without changing the verdict on the cases tested, and 11/12 (92\%)
correct on the curated development set. The same results also show this
headline figure does not generalize on its own: on an unfiltered random
sample, precision held near the curated figure on one backend (92.9\%,
\texttt{iec}) but collapsed on another (60.3\%, \texttt{fit}/Qwen), showing
that the accuracy story is backend-dependent as much as it is
dataset-dependent.

\section{Contributions and Significance}

This thesis's contributions are scoped honestly to what was actually built,
not to the fuller cross-candidate comparison originally targeted
(Chapter~\ref{Chapter4}). As of this writing, one
candidate has a working, measured implementation; the significance below
follows from that one candidate, not from a comparison across methods.

\textbf{A measured, not asserted, cost-efficiency result.} The literature
surveyed in Chapter~\ref{Chapter3} identifies cost as a persistent, poorly
standardized concern across LLM-based detectors (RQ2, RQ3) but rarely
isolates which specific design choice buys which specific saving. This
thesis contributes a design history (Table~\ref{tab:design-history}) that
attributes each cost reduction to a specific, independently-tested
mechanism -- caching versus connection-verified early exit -- with paired
before/after numbers for each, rather than a single aggregate saving
figure whose source is hard to attribute.

\textbf{Evidence that cost-saving and correctness-preserving are separate
engineering problems.} The same design history shows, with dated,
root-caused examples, that a cost-reduction mechanism which appears to
work (naive early exit; same-class connection without verification) can
silently change a verdict, and that closing each gap was a distinct
engineering step, not a byproduct of the original optimization. This is a
concrete, reproducible instance of a pattern the surveyed literature
discusses more abstractly (Chapter~\ref{Chapter3}, RQ3's discussion of
verification mechanisms being themselves susceptible to error).

\textbf{Empirical confirmation, on a fresh sample, of two gaps the SLR
identified from the literature alone.} The unfiltered 750-APK comparison
(Section~\ref{sec:threats-ch5}) is independent evidence -- not a citation
of prior work -- for two claims Chapter~\ref{Chapter3} makes from reading
other papers: that reliance on static analysis alone leaves a persistent
blind spot (confirmed here by a real missed detection carried entirely in
a native \texttt{.so} library, Chapter~\ref{Chapter5}), and that
verification mechanisms do not catch every reliability failure (confirmed
here by false positives that are structurally accurate but semantically
misjudged, evading LAMD's own Factual Consistency Verification).

\textbf{A reusable evaluation framework.} Independent of any single
candidate's results, Chapter~\ref{Chapter4}'s dataset, metric, and
protocol definitions are built to accept additional candidates from the
Chapter~\ref{Chapter3} taxonomy without redefinition, which is the
practical prerequisite for the cross-candidate comparison this thesis did
not have time to complete.

\section{Limitations}\label{sec:limitations-ch6}

The measurement-level limitations of the results already reported --
sample size, LLM non-determinism, an unswept cache threshold -- are
discussed in Section~\ref{sec:threats-ch5} and are not repeated here. This
section instead states the boundaries of what the thesis attempted at
all.

\textbf{Scope: one candidate, not a comparison.} The thesis's own stated
goal (Chapter~\ref{Chapter4}) was a common framework for comparing
multiple LLM-based candidates. Only one candidate reached a working,
measured implementation. Every accuracy, cost, and reliability finding in
this thesis is therefore a finding about the LAMD-based design
specifically, not a general claim about LLM-based Android malware
detection as a class.

\textbf{The 750-APK evaluation is unfinished.} The random-sample
comparison central to Chapter~\ref{Chapter5}'s Discussion and Threats to
Validity was, at the time of writing, an in-progress run (596/750 and
215/750 complete on the two backends), and the \texttt{fit} backend's
figures mix pre- and post-rebalance results from a dataset composition
change made partway through. The directional finding -- precision is
backend-dependent -- is unlikely to reverse, but the exact percentages
reported are not final numbers.

\textbf{The RAG-grounding exploration is not a second candidate.} A
retrieval-augmented prompt-grounding technique was implemented and tested
on top of the LAMD-based candidate (Chapter~\ref{Chapter4}), but it was
never evaluated under the common framework's accuracy/cost metrics on more
than a single APK, and it reuses the LAMD-based candidate's own pipeline
rather than constituting an independent detection architecture. It is
reported as an internal exploration, not as evidence toward a second,
independently-evaluated candidate.

\section{Future Work}

\textbf{Complete the in-progress 750-APK evaluation.} The clearest
near-term action is finishing the random-sample run on both backends and
recomputing Section~\ref{sec:threats-ch5}'s figures on the full,
post-rebalance 750 APKs, since the current numbers are explicitly
provisional.

\textbf{Address Tier-3 non-determinism.} A self-consistency or
majority-vote Tier-3 (issuing the final-verdict call multiple times and
taking the majority) was discussed as a candidate fix during development
(Section~\ref{sec:threats-ch5}) but not implemented, since it
proportionally increases per-APK cost -- itself a concrete instance of the
cost/reliability trade-off this thesis's cost-efficiency framing puts
front and center, and a natural next design-history entry to measure
rather than assume.

\textbf{Sweep the cache similarity threshold.} The $\geq 0.88$ cosine
threshold was fixed by inspection; a systematic sweep would show whether
cache-hit rate, cost, and downstream accuracy trade off favorably at other
settings, closing the gap noted in Section~\ref{sec:threats-ch5}.

\textbf{Implement additional candidates to complete the cross-candidate
comparison.} This is the most direct way to close the scope gap in
Section~\ref{sec:limitations-ch6} above: at least one RAG-based and one
fine-tuning-based candidate, representative of the two families identified
in Chapter~\ref{Chapter3}'s taxonomy (RQ1) but not yet implemented here,
would let Chapter~\ref{Chapter4}'s evaluation framework fulfill its
original purpose.

\textbf{Evaluate a cascaded triage architecture.} Chapter~\ref{Chapter5}'s
Discussion motivates this concretely, not abstractly: the false-positive
pattern observed on the \texttt{fit} backend (ordinary login flows,
common bundled-library reflection patterns) is the kind of case a
lightweight pre-filter could plausibly screen out before full LLM
reasoning is invoked, in the spirit of Chapter~\ref{Chapter3}'s RQ4
recommendation, and this thesis's own false-positive data is a ready-made
test set for that idea.
